
\documentclass[11pt,a4paper]{article}

\usepackage[a4paper,margin=25mm]{geometry}
\usepackage[T1]{fontenc}
\usepackage[utf8]{inputenc}

\pagestyle{empty}

\begin{document}

\section*{Statement of Significance}

Mean-strain distortion sits between two familiar limits: rapid-distortion
theory, which freezes nonlinear relaxation, and Reynolds-stress modelling, which
removes the spectrum. This paper closes that gap by embedding production, rapid
pressure--strain redistribution and scale compression into ODT while retaining
an evolving spectrum. The work is fully developed through analysis,
implementation checks, RDT verification and DNS/experiment comparisons. The main
fluid-mechanical insight is that at $\chi=O(1)$ strain and eddy relaxation act
together, producing a broadband distorted spectrum rather than a linearly
translated one.

\end{document}
